% ============================================================
% M3 S3 练习件·W1 臂A —— 课时03 方向向量与法向量（2.2.1下）＝照抄母版件（课时01）体例
% 生成：M3 成卷轮 S3 W1 臂A｜2026-09-14｜编译 cwd＝本目录（中文路径禁 \input 绝对路径）
%   xelatex main-true.tex ×2 → 含详解印本；xelatex main-false.tex ×2 → 纯题版（单源双本）
% 输入链（全只读）：题面库 课时03-方向向量与法向量.md（题面逐字装配）＋ -答案侧.md（值逐字回嵌）
%   ＋ 定稿/批A-课时03-方向向量与法向量.md（详解回嵌源，0913命制入槽 E12~E16 并指件内
%   【亲算落账】格——亲算落账系过程账，不入印面）；toolchain＝qp-m3.sty（钉后版
%   md5 7c3930362be8a0a2bdf21bbf8ac16573，本地挂载副本，破损即停）。
% 母版体例：详见《练习件母版体例说明.md》（课时01 件）；门谱读数见 过程件 工作区/_tmpM3S3W1臂A0914/。
% 键账：件内 ansblock 键集 ≡ 题面库片练键集 ≡ manifest 练键序 E1~E16（16 键，零缺漏零多余）；
%   拓区隔离：2章-拓/测/滚 键不入本件（S4 拓展册收；本片拓展槽空置无 T 键，定稿§三/§四 在案）；
%   导学键 G1~G5 属导学件域，亦不入。
% 卷式例外案B：练习件不属卷式——答案块物理落点恒为题后 ansblock，禁卷末附卷制；
%   全线括线模（导言 \ansblockgrayfalse 一次置定，件内不切换；灰底盒不可跨栏断，练习件禁用）。
% 题序＝manifest 键序 E1~E16（零重排）；印面题号 1~16 连号（\tihao 号＝\ansitem 号同键）；
%   三档切位 1~7／8~14／15~16（照库序切位，非难度重排；E12~E16＝0913命制入槽新命对，
%   题面照录轮4回修件全文一字未动）。
% 槽配实录：单选1（E12）＋多选0（manifest 期望值 0）＋填空3（E1/E13/E16）＋解答12（E2~E11/E14/E15）；
%   10选4填2解 为波次方案基线槽配，本片实配照题面侧头标（批A §四满足度表同口径：选择缺9、
%   填空缺1、解答超10 形态偏差按实登记，缺5已由 0913命制入槽闭合）。
% 选项槽位：默认四连排（0.25\linewidth−0.25\lxhang），超宽降档梯 四连排→两连排
%   （0.5\linewidth−0.5\lxhang−0.25em）→单列 \lxopt；禁缩字号/负 kern/删标点凑宽。
%   本片实测降档：见值台账「降档登记」＋回执（槽宽门真算读数）。
% 解答书写区：无小问默认 16mm（E3/E4/E5/E8/E11/E15），两问 24mm（E2/E6/E7/E9/E10/E14，
%   E4 两问＝「A，B，C共线吗？A，B，D呢？」双问句）；纯题档吞 ansblock 不吞 \liubai。
% 填空印答：E1 双空 \kongda{2}＋\kongda{4/3}、E13 \kongda{2x−y−5=0}、E16 \kongda{−√3/3}
%   （详解版印答、纯题版退定宽空线，两档等形）。
% 难度词：照题面侧头标第 4 字段简/中档词逐字；知识点 N 系预排（本课时知识点一＝直线的
%   方向向量、二＝直线的法向量，照教材结构-第2章 2.2.1 条目 2/3 两知识点切分；导学件课时03
%   尚未落盘，落盘后如异动按变更协议回核）。
% ============================================================
\documentclass[fontset=none]{ctexart}
\usepackage{qp-m3}
% —— 印面逐字符号路由补钉（承导学件母版；− 走 TNR、▱ NSC 子块；禁改 sty 保 md5） ——
\xeCJKDeclareCharClass{Default}{"2212}%
\setCJKfamilyfont{nscblk}[Path=C:/提示词/工作区/字替对照-0909/variantF/fonts/,BoldFont=NSC-w600.ttf]{NSC-w500.ttf}%
\xeCJKDeclareSubCJKBlock{nscblk}{"25B1}%
\newcommand{\pxparallelogram}{{\CJKfamily{nscblk}▱}}%
% —— 断行微调（承导学件母版实测档，三零门承重；\emergencystretch 不另设） ——
\xeCJKsetup{CJKglue={\hskip 0.10em plus 0.08em minus 0.05em}}%
% —— 练习件全线括线模（S3 波次方案 §四.3：导言一次置定、件内不切换） ——
\ansblockgrayfalse
% —— 页脚件名（qp-headfoot 复用入口） ——
\renewcommand{\qpzhangming}{第二章\quad 平面解析几何}%
\renewcommand{\qpjianming}{练习件}%

\begin{document}
%装配轮S6三本跨本连续页码（G7方案B）setcounter 注入位
\setcounter{page}{168}

\keshi{第3课时\quad 方向向量与法向量}
\begin{multicols}{2}
\emergencystretch=1em

% ============================================================
% 基础巩固（题1~7＝E1~E7：填空1＋解答6；题后紧跟答案制，全线括线 ansblock）
% ============================================================
\dangbadge{基}{础}{巩}{固}

\tihao{1}\zu{直线的方向向量与法向量×填空} \tieside{简单(知识点一)}已知直线\(l\)经过点\(P(3,m)\)和点\(Q(m,-2)\)，直线\(l\)的方向向量为\((2,4)\)，则直线\(l\)的斜率为\kongda{2}，实数\(m\)的值为\kongda{4/3}．

\begin{ansblock}[2章-练-课时03-E1]
% ans:2章-练-课时03-E1
\ansitem{1}{2；4/3}
\ansnote{详解}{由方向向量\((2,4)\)得直线\(l\)的斜率 \(k=4/2=2\)．\par\noindent 由两点斜率公式 \(k=(m-(-2))/(3-m)=(m+2)/(3-m)\)（\(m\neq 3\)），令\((m+2)/(3-m)=2\)，得\(m+2=6-2m\)，\(3m=4\)，\(m=4/3\)．\par\noindent 故斜率为2，\(m=4/3\)．}
\end{ansblock}

\tihao{2}\zu{直线的方向向量与法向量×解答} \tieside{简单(知识点一)}（1）设坐标平面内三点\(A(m,-m-3)\)，\(B(2,m-1)\)，\(C(-1,4)\)，若直线\(AC\)的斜率等于直线\(BC\)的斜率的3倍，求实数\(m\)的值；（2）已知直线\(l_1\)的方向向量为\(\overrightarrow{n}=(2,1)\)，直线\(l_2\)的倾斜角是直线\(l_1\)倾斜角的2倍，求直线\(l_2\)的斜率．
\liubai[24mm]{解答书写区}

\begin{ansblock}[2章-练-课时03-E2]
% ans:2章-练-课时03-E2
\ansitem{2}{(1) 1或2；(2) 4/3}
\ansnote{详解}{(1) \(k_{AC}=[4-(-m-3)]/(-1-m)=(m+7)/(-m-1)\)（\(m\neq -1\)），\(k_{BC}=(4-m+1)/(-1-2)=(5-m)/(-3)\)．\par\noindent 由\(k_{AC}=3k_{BC}\)：\((m+7)/(-m-1)=3\cdot(5-m)/(-3)=-(5-m)=m-5\)，\par\noindent 去分母得 \(m+7=(m-5)(-m-1)=-m^{2}+4m+5\)，即 \(m^{2}-3m+2=0\)，解得\(m=1\)或\(m=2\)，经检验均使两斜率存在，符合题意．\par\noindent (2) 设\(l_1\)的倾斜角为\(\alpha\)，则\(l_2\)的倾斜角为\(2\alpha\)．由方向向量\(\overrightarrow{n}=(2,1)\)得 \(\tan\alpha=1/2\)，\par\noindent 故\(l_2\)的斜率 \(\tan 2\alpha=2\tan\alpha/(1-\tan^{2}\alpha)=1/(1-1/4)=4/3\)．}
\end{ansblock}

\tihao{3}\zu{直线的方向向量与法向量×解答} \tieside{简单(知识点一)}已知直线\(l\)经过点\(A(-2,0)\)与\(B(-5,3)\)，求直线\(l\)的一个方向向量、斜率\(k\)与倾斜角\(\theta\)．
\liubai[16mm]{解答书写区}

\begin{ansblock}[2章-练-课时03-E3]
% ans:2章-练-课时03-E3
\ansitem{3}{方向向量如(1,−1)（答案不唯一）；k＝−1；θ＝135°}
\ansnote{详解}{\(\overrightarrow{AB}=(-5+2,\ 3-0)=(-3,3)\)是直线\(l\)的一个方向向量（与之共线的非零向量如\((1,-1)\)也可）．\par\noindent 斜率 \(k=(3-0)/(-5-(-2))=3/(-3)=-1\)（也可由方向向量\((1,-1)\)得\(k=-1/1=-1\)）．\par\noindent 由\(\tan\theta=-1\)且\(\theta\in[0^\circ,180^\circ)\)，得\(\theta=135^\circ\)．}
\end{ansblock}

\tihao{4}\zu{直线的方向向量与法向量×解答} \tieside{简单(知识点一)}已知\(A(-1,-3)\)，\(B(0,-1)\)，\(C(1,2)\)，\(D(3,5)\)，则\(A\)，\(B\)，\(C\)共线吗？\(A\)，\(B\)，\(D\)呢？
\liubai[24mm]{解答书写区}

\begin{ansblock}[2章-练-课时03-E4]
% ans:2章-练-课时03-E4
\ansitem{4}{A、B、C不共线；A、B、D共线．}
\ansnote{详解}{\(\overrightarrow{AB}=(0+1,-1+3)=(1,2)\)，\(\overrightarrow{AC}=(1+1,2+3)=(2,5)\)．\par\noindent 因为\(1\times 5-2\times 2=1\neq 0\)，\(\overrightarrow{AB}\)与\(\overrightarrow{AC}\)不共线，且两向量有公共起点\(A\)，所以\(A\)、\(B\)、\(C\)不共线．\par\noindent \(\overrightarrow{AD}=(3+1,5+3)=(4,8)=4\overrightarrow{AB}\)，\(\overrightarrow{AB}\)与\(\overrightarrow{AD}\)共线且有公共起点\(A\)，所以\(A\)、\(B\)、\(D\)共线．}
\end{ansblock}

\tihao{5}\zu{直线的方向向量与法向量×解答} \tieside{简单(知识点一)}判断命题“如果\(A\)，\(B\)，\(C\)是平面直角坐标系中的三个不同的点，则这三点共线的充要条件是\(\overrightarrow{AB}\)与\(\overrightarrow{BC}\)共线”的真假．
\liubai[16mm]{解答书写区}

\begin{ansblock}[2章-练-课时03-E5]
% ans:2章-练-课时03-E5
\ansitem{5}{真命题．证明见详解．}
\ansnote{详解}{必要性：若\(A\)、\(B\)、\(C\)共线，则直线上的向量\(\overrightarrow{AB}\)、\(\overrightarrow{BC}\)及与直线平行的向量都是该直线的方向向量，故\(\overrightarrow{AB}\)与\(\overrightarrow{BC}\)共线．\par\noindent 充分性：若\(\overrightarrow{AB}\)与\(\overrightarrow{BC}\)共线，则直线\(AB\)与直线\(BC\)的方向相同或相反；又两直线过公共点\(B\)，由过一点且方向确定的直线唯一，直线\(AB\)与直线\(BC\)是同一条直线，故\(A\)、\(B\)、\(C\)三点共线．\par\noindent 综上命题为真．}
\end{ansblock}

\tihao{6}\zu{直线的方向向量与法向量×解答} \tieside{中档(知识点一)}(1) 如果直线\(l\)过点\(P(-1,-2)\)，且直线\(l\)的方向向量为\(a=(1,-2)\)，求直线\(l\)的方程；(2) 如果直线\(l\)过点\(P(x_0,y_0)\)，且直线\(l\)的方向向量为\(a=(u,v)\)，求直线\(l\)的方程．
\liubai[24mm]{解答书写区}

\begin{ansblock}[2章-练-课时03-E6]
% ans:2章-练-课时03-E6
\ansitem{6}{(1) y＝−2x−4（即2x＋y＋4＝0）；(2) 当u≠0时为y−y₀＝(v/u)(x−x₀)；当u＝0时为x＝x₀．}
\ansnote{详解}{(1) 方向向量\(a=(1,-2)\)的横坐标\(1\neq 0\)，故斜率\(k=-2/1=-2\)，由点斜式得 \(y-(-2)=-2[x-(-1)]\)，即 \(y=-2x-4\)．\par\noindent (2) 设\(Q(x,y)\)是直线\(l\)上任意一点，则\(\overrightarrow{PQ}=(x-x_0,\ y-y_0)\)与方向向量\(a=(u,v)\)共线，得 \(v(x-x_0)-u(y-y_0)=0\)．\par\noindent 当\(u\neq 0\)时化为 \(y-y_0=(v/u)(x-x_0)\)；当\(u=0\)时（此时\(v\neq 0\)）得 \(x-x_0=0\)，即\(x=x_0\)（斜率不存在的竖直线）．}
\end{ansblock}

\tihao{7}\zu{直线的方向向量与法向量×解答} \tieside{中档(知识点二)}(1) 如果直线\(l\)过点\(P(1,3)\)，且直线\(l\)的法向量为\(a=(-3,1)\)，求直线\(l\)的方程；(2) 如果直线\(l\)过点\(P(x_0,y_0)\)，且直线\(l\)的法向量为\(a=(u,v)\)，求直线\(l\)的方程．
\liubai[24mm]{解答书写区}

\begin{ansblock}[2章-练-课时03-E7]
% ans:2章-练-课时03-E7
\ansitem{7}{(1) y＝3x；(2) u(x−x₀)＋v(y−y₀)＝0，即ux＋vy−ux₀−vy₀＝0．}
\ansnote{详解}{(1) 设\(Q(x,y)\)是直线\(l\)上任意一点，则\(\overrightarrow{PQ}=(x-1,\ y-3)\)与法向量\(a=(-3,1)\)垂直，\(a\cdot\overrightarrow{PQ}=0\)，\par\noindent 即 \(-3(x-1)+1\cdot(y-3)=0\)，化简得 \(y=3x\)．\par\noindent (2) 同法：\(a\cdot\overrightarrow{PQ}=u(x-x_0)+v(y-y_0)=0\)（\(u\)、\(v\)不全为0，因法向量是非零向量）．这就是直线的点法式方程；\par\noindent 当\(v\neq 0\)时可化为斜截形 \(y=-(u/v)x+(ux_0+vy_0)/v\)，斜率\(k=-u/v\)（与“方向向量\((v,-u)\)读斜率\(-u/v\)”一致）．}
\end{ansblock}

% ============================================================
% 综合提升（题8~14＝E8~E14：解答6（E8~E11/E14）＋单选1（E12）＋填空1（E13））
% ============================================================
\dangbadge{综}{合}{提}{升}

\tihao{8}\zu{直线的方向向量与法向量×解答} \tieside{简单(知识点二)}向量\((x_0,y_0)\)与\((-y_0,x_0)\)中，如果其中一个为直线\(l\)的一个方向向量，则另一个一定是直线\(l\)的一个法向量吗？
\liubai[16mm]{解答书写区}

\begin{ansblock}[2章-练-课时03-E8]
% ans:2章-练-课时03-E8
\ansitem{8}{一定．理由见详解．}
\ansnote{详解}{设\((x_0,y_0)\)是\(l\)的一个方向向量．由方向向量非零知\((x_0,y_0)\neq(0,0)\)，从而\((-y_0,x_0)\neq(0,0)\)．\par\noindent 又 \((x_0,y_0)\cdot(-y_0,x_0)=-x_0y_0+y_0x_0=0\)，即\((-y_0,x_0)\)与\(l\)的方向向量垂直，于是以\((-y_0,x_0)\)为方向向量的有向线段所在直线与\(l\)垂直，\par\noindent 由法向量的定义，\((-y_0,x_0)\)是\(l\)的一个法向量．\par\noindent 反之，若\((-y_0,x_0)\)为\(l\)的方向向量，同理\((0,0)\neq(-y_0,x_0)\)\(\Longleftrightarrow\)\((x_0,y_0)\neq(0,0)\)，且两向量数量积为0，故\((x_0,y_0)\)一定是\(l\)的法向量．\par\noindent 所以结论成立（教材正文中“特别地，当\(x_0\)与\(y_0\)不全为0时…”的前提恰好被“方向向量非零”自动满足）．}
\end{ansblock}

\tihao{9}\zu{直线的方向向量与法向量×解答} \tieside{中档(知识点一)}已知\(P\)是直线\(l\)上一点，且\(n\)是直线\(l\)的一个方向向量，分别根据下列条件求直线\(l\)的方程：(1) \(P(3,-5)\)，\(n=(1,2)\)；(2) \(P(0,5)\)，\(n=(3,-4)\)．
\liubai[24mm]{解答书写区}

\begin{ansblock}[2章-练-课时03-E9]
% ans:2章-练-课时03-E9
\ansitem{9}{(1) y＝2x−11；(2) y＝−(4/3)x＋5（即4x＋3y−15＝0）}
\ansnote{详解}{(1) \(n=(1,2)\)横坐标不为0，\(k=2/1=2\)，由点斜式 \(y+5=2(x-3)\)，即 \(y=2x-11\)．\par\noindent (2) \(n=(3,-4)\)横坐标不为0，\(k=-4/3\)，由点斜式 \(y-5=-(4/3)(x-0)\)，即 \(y=-(4/3)x+5\)．}
\end{ansblock}

\tihao{10}\zu{直线的方向向量与法向量×解答} \tieside{中档(知识点二)}已知\(P\)是直线\(l\)上一点，且\(v\)是直线\(l\)的一个法向量，分别根据下列条件求直线\(l\)的方程：(1) \(P(1,2)\)，\(v=(3,-4)\)；(2) \(P(-1,2)\)，\(v=(3,4)\)．
\liubai[24mm]{解答书写区}

\begin{ansblock}[2章-练-课时03-E10]
% ans:2章-练-课时03-E10
\ansitem{10}{(1) 3x−4y＋5＝0；(2) 3x＋4y−5＝0}
\ansnote{详解}{(1) 设\(Q(x,y)\)在\(l\)上，\(v\cdot\overrightarrow{PQ}=3(x-1)-4(y-2)=0\)，化简得 \(3x-4y+5=0\)．\par\noindent (2) 同法：\(3(x+1)+4(y-2)=0\)，化简得 \(3x+4y-5=0\)．}
\end{ansblock}

\tihao{11}\zu{直线的方向向量与法向量×解答} \tieside{中档(知识点二)}已知三点\(A(1,4)\)，\(B(3,3)\)，\(C(4,5)\)．求证：\(AB\perp BC\)（要求用直线的方向向量与法向量的语言完成）．
\liubai[16mm]{解答书写区}

\begin{ansblock}[2章-练-课时03-E11]
% ans:2章-练-课时03-E11
\ansitem{11}{证明见详解．}
\ansnote{详解}{\(\overrightarrow{AB}=(3-1,3-4)=(2,-1)\)，故直线\(AB\)的一个方向向量为\(a=(2,-1)\)；\(\overrightarrow{BC}=(4-3,5-3)=(1,2)\)．\par\noindent 取\(v=(1,2)\)，验证\(v\)与\(a\)垂直：\(v\cdot a=1\times 2+2\times(-1)=0\)，且\(v\neq 0\)，所以\(v\)是直线\(AB\)的一个法向量．\par\noindent 而\(\overrightarrow{BC}=v\)，即直线\(BC\)的方向向量就是直线\(AB\)的法向量，由法向量的定义，直线\(BC\)与直线\(AB\)垂直，故\(AB\perp BC\)．}
\end{ansblock}

\tihao{12}\zu{直线的方向向量与法向量×单选} \tieside{简单(知识点一)}直线\(l\)经过\(A(1,2)\)、\(B(3,6)\)两点，则\(l\)的一个方向向量可以是\nobreak(\kongwei)
\bindopt {\fontsize{10.5pt}{\qpTJgLeadLiti}\selectfont \makebox[\dimexpr0.25\linewidth-0.25\lxhang\relax][l]{A．\((2,1)\)}\makebox[\dimexpr0.25\linewidth-0.25\lxhang\relax][l]{B．\((1,2)\)}\makebox[\dimexpr0.25\linewidth-0.25\lxhang\relax][l]{C．\((-2,1)\)}\makebox[\dimexpr0.25\linewidth-0.25\lxhang\relax][l]{D．\((1,-2)\)}\par}

\begin{ansblock}[2章-练-课时03-E12]
% ans:2章-练-课时03-E12
\ansitem{12}{B}
\ansnote{详解}{\(AB=(3-1,\ 6-2)=(2,4)=2(1,2)\)，故 \((1,2)\) 为方向向量，选 B。}
\end{ansblock}

\tihao{13}\zu{直线的方向向量与法向量×填空} \tieside{简单(知识点二)}已知直线\(l\)的一个法向量为\(n=(2,-1)\)，且\(l\)过点\(P(3,1)\)，则\(l\)的方程为\kongda{2x−y−5=0}。

\begin{ansblock}[2章-练-课时03-E13]
% ans:2章-练-课时03-E13
\ansitem{13}{2x−y−5=0}
\ansnote{详解}{法向量 \((2,-1)\)，故 \(l\) 具形式 \(2x-y+C=0\)。代 \(P(3,1)\)：\(6-1+C=0\)，即 \(C=-5\)。故 \(2x-y-5=0\)。}
\end{ansblock}

\tihao{14}\zu{直线的方向向量与法向量×解答} \tieside{中档(知识点二)}已知直线\(l_1\): \(3x-y+1=0\)与\(l_2\): \(ax+2y-6=0\)。(1) 写出\(l_1\)的一个法向量；(2) 若\(l_1\perp l_2\)，求\(a\)的值，并用法向量语言说明理由。
\liubai[24mm]{解答书写区}

\begin{ansblock}[2章-练-课时03-E14]
% ans:2章-练-课时03-E14
\ansitem{14}{(1) n₁=(3,−1)（任一非零倍数均可）　(2) a=2/3}
\ansnote{详解}{(1) 由一般式系数直接读出 \(n_1=(3,-1)\)。(2) \(n_2=(a,2)\)。同一平面内每条直线的法向量与其方向向量互相垂直，将两直线绕任一点同旋 \(90^\circ\)，方向向量恰转到各自法向量方向且夹角不变，\par\noindent 故 \(l_1\perp l_2\)\(\Longleftrightarrow\)\(n_1\perp n_2\)\(\Longleftrightarrow\)\(n_1\cdot n_2=0\)\(\Longleftrightarrow\)\(3a+(-1)\times 2=0\)\(\Longleftrightarrow\)\(a=2/3\)。}
\end{ansblock}

% ============================================================
% 思维探索（题15~16＝E15~E16：解答1＋填空1，居末档照库序；两题均系 0913命制入槽新命对）
% ============================================================
\dangbadge{思}{维}{探}{索}

\tihao{15}\zu{直线的方向向量与法向量×解答} \tieside{中档(知识点一)}已知直线\(l\): \((m+1)x+2y-6=0\)的一个方向向量为\(v=(2,\ m-1)\)，求实数\(m\)的值。
\liubai[16mm]{解答书写区}

\begin{ansblock}[2章-练-课时03-E15]
% ans:2章-练-课时03-E15
\ansitem{15}{m=0}
\ansnote{详解}{直线 \(Ax+By+C=0\) 的方向向量可取 \((B,\ -A)\)，故 \(l\) 的方向向量 \(t=(2,\ -(m+1))\)。\par\noindent \(t\parallel v\)（横坐标已同为 2），即 \(-(m+1)=m-1\)，\par\noindent 解得 \(2m=0\)，故 \(m=0\)。}
\end{ansblock}

\tihao{16}\zu{直线的方向向量与法向量×填空} \tieside{简单(知识点一)}直线\(l\)过点\(P(0,1)\)，且一个方向向量为\(v=(1,\sqrt{3})\)，则\(l\)在\(x\)轴上的截距为\kongda{−√3/3}。

\begin{ansblock}[2章-练-课时03-E16]
% ans:2章-练-课时03-E16
\ansitem{16}{−√3/3}
\ansnote{详解}{斜率 \(k=\sqrt{3}/1=\sqrt{3}\)，\(l\): \(y=\sqrt{3}x+1\)。令 \(y=0\)：\(x=-1/\sqrt{3}=-\sqrt{3}/3\)。}
\end{ansblock}

% ---- 尾块（\tailfill 置 \end{multicols} 前；无参＝内容尾随框，每件恰 1 处） ----
\tailfill

\end{multicols}
\end{document}
